%-------------------------
% Resume in Latex
% Based on Jake's Resume Template (https://github.com/jakegut/resume / sourabhbajaj)
%
% ATS notes:
%  - Single column, standard section headings (Experience / Education &
%    Certifications / Technical Skills / Projects), no text boxes, no images,
%    no multi-column layout. Selectable text with \pdfgentounicode=1 so
%    parsers read it cleanly.
%  - Link text uses ulem's \uline, not \underline: \underline typesets in math
%    mode, which ignores \small and extracts worse. Do not switch it back.
%  - Every claim below is backed by code that exists in a repo. Do not re-add
%    a bullet you cannot open a file for.
%------------------------

\documentclass[letterpaper,11pt]{article}

\usepackage{latexsym}
\usepackage[empty]{fullpage}
\usepackage{titlesec}
\usepackage{marvosym}
\usepackage[usenames,dvipsnames]{color}
\usepackage{verbatim}
\usepackage{enumitem}
\usepackage[hidelinks]{hyperref}
\usepackage{fancyhdr}
\usepackage[english]{babel}
\usepackage[normalem]{ulem}
\usepackage{tabularx}
\input{glyphtounicode}

\pagestyle{fancy}
\fancyhf{}
\fancyfoot{}
\renewcommand{\headrulewidth}{0pt}
\renewcommand{\footrulewidth}{0pt}

% Adjust margins
% Text block is 7.9in x 10.55in on letter: 0.30in side margins, 0.30in top,
% 0.15in bottom. This is the floor -- 0.15in bottom is inside the
% unprintable region of many printers, so the PDF is screen/ATS-first.
% Reclaim space here rather than from itemsep or heading gaps, both of
% which cause visible collisions (see the notes on those commands).
\addtolength{\oddsidemargin}{-0.7in}
\addtolength{\evensidemargin}{-0.7in}
\addtolength{\textwidth}{1.4in}
\addtolength{\topmargin}{-0.7in}
\addtolength{\textheight}{1.55in}

\urlstyle{same}

\raggedbottom
\raggedright
\setlength{\tabcolsep}{0in}

% Sections formatting
\titleformat{\section}{
  \vspace{-8pt}\scshape\raggedright\large
}{}{0em}{}[\color{black}\titlerule \vspace{-8pt}]

% Ensure that generated pdf is machine readable / ATS parsable
\pdfgentounicode=1

\newcommand{\graddate}{2023 -- 2027}

\newcommand{\introle}{Machine Learning Intern}
\newcommand{\intdates}{2026 -- Present}

%-------------------------
% Custom commands
% No negative \vspace here: it used to shrink the gap *between* bullets while
% the leading *inside* a wrapped bullet stayed at full \baselineskip, so a
% bullet's last line sat closer to the next bullet than to its own first line.
% Inter-item space must stay larger than intra-item leading -- see itemsep below.
\newcommand{\resumeItem}[1]{
  \item\small{#1}
}

\newcommand{\resumeSubheading}[4]{
  \vspace{-2pt}\item
    \begin{tabular*}{0.97\textwidth}[t]{l@{\extracolsep{\fill}}r}
      \textbf{#1} & #2 \\
      \textit{\small#3} & \textit{\small #4} \\
    \end{tabular*}\vspace{-8pt}
}

\newcommand{\resumeProjectHeading}[2]{
    \item
    \begin{tabular*}{0.97\textwidth}{l@{\extracolsep{\fill}}r}
      \small#1 & \small#2 \\
    \end{tabular*}\vspace{-5pt}
}

\newcommand{\resumeSubItem}[1]{\resumeItem{#1}\vspace{-4pt}}

\renewcommand\labelitemii{$\vcenter{\hbox{\tiny$\bullet$}}$}

\newcommand{\resumeSubHeadingListStart}{\begin{itemize}[leftmargin=0.15in, label={}, itemsep=0pt]}
% Space is reclaimed at the boundaries next to bold headings, never from
% itemsep -- a heading separates blocks visually on its own, so tightening
% there costs nothing, while tightening itemsep reintroduces the bug above.
\newcommand{\resumeSubHeadingListEnd}{\end{itemize}\vspace{-11pt}}
\newcommand{\resumeItemListStart}{\begin{itemize}[itemsep=2pt, parsep=0pt, topsep=0pt]}
% -6pt, not lower: a \resumeProjectHeading follows this in the Projects
% section, and anything near -11pt drives the heading into the last bullet.
\newcommand{\resumeItemListEnd}{\end{itemize}\vspace{-6pt}}

%-------------------------------------------
%%%%%%  RESUME STARTS HERE  %%%%%%%%%%%%%%%%%%%%%%%%%%%%%%%%%%%%%%%%

\begin{document}

%----------HEADING----------
\begin{center}
    \textbf{\Huge \scshape Divyansh Shukla} \\ \vspace{2pt}
    % \scriptsize, not \small: the single-line contact row measures 635pt at
    % \small against a 571pt text block -- it cannot fit at any margin. See
    % the measurement note in the commit message.
    \scriptsize
    \href{mailto:divyanshmohanshukla@hotmail.com}{\uline{divyanshmohanshukla@hotmail.com}} $|$
    +91 7827910144 $|$
    \href{https://github.com/divyanshuklai}{\uline{github.com/divyanshuklai}} $|$
    \href{https://www.linkedin.com/in/divyansh-shukla-ai/}{\uline{linkedin.com/in/divyansh-shukla-ai}} $|$
    \href{https://divyanshuklai.github.io}{\uline{divyanshuklai.github.io}}
\end{center}

%-----------EDUCATION-----------
\section{Education}
  \resumeSubHeadingListStart
    \resumeSubheading
      {Vellore Institute of Technology}{\graddate}
      {B.Tech, Computer Science and Engineering (Artificial Intelligence \& Machine Learning)}{CGPA: 9.03 / 10}
  \resumeSubHeadingListEnd

%-----------EXPERIENCE-----------
\section{Experience}
  \resumeSubHeadingListStart
    \resumeSubheading
      {Power Pulse India}{\intdates}
      {\introle}{India}
      \resumeItemListStart
        \resumeItem{Cut missed pedestrians 26\% and raised mAP50-95 from 0.360 to 0.420 for a driver-assistance thermal detector, training a deployment-aware YOLOv8n on FLIR ADAS v2 across 6 road-safety classes in 100 epochs against 150, and trading 3 points of precision for the recall that safety case requires.}
        \resumeItem{Identified annotation inconsistency and input resolution as the primary accuracy bottlenecks through per-class crop audits and COCO size-bucket analysis: 82\% of pedestrian labels fall under 32x32 px, and labelled trucks were delivery vans also annotated car and bus, with 37\% of true trucks drawing no prediction at all.}
        \resumeItem{Ran a four-experiment resolution and architecture sweep to cut deployment cost, eliminating a P2 detection head that scored below plain 640px at 1.51x the compute and 4x the detection candidates, and identifying 800px as capturing 85\% of the small-object recall gain at 1.56x compute against 960px's 2.25x.}
      \resumeItemListEnd
  \resumeSubHeadingListEnd

%-----------CERTIFICATIONS-----------
\section{Certifications}
 \begin{itemize}[leftmargin=0.15in, label={}]
    \small{\item{
     \textbf{NPTEL}{: Applied Accelerated AI, Deep Learning with Computer Vision} \hspace{3pt}$|$\hspace{3pt} \textbf{Coursera}{: Deep Learning Specialization} \hspace{3pt}$|$\hspace{3pt} \textbf{AWS}{: AWS Cloud Architecting}
    }}
 \end{itemize}

%-----------SKILLS-----------
\section{Technical Skills}
 \begin{itemize}[leftmargin=0.15in, label={}]
    \small{\item{
     \textbf{Languages}{: Python, SQL, C++, CUDA} \\
     \textbf{ML Frameworks}{: PyTorch, PyTorch Lightning, Hugging Face Transformers, TRL, Unsloth, Ultralytics (YOLO), Stable-Baselines3, vLLM, llama.cpp} \\
     \textbf{Deep Learning}{: flow/diffusion models, transformers, YOLO, generation, inpainting, detection, visual question answering (VLM), reinforcement learning (PPO)} \\
     \textbf{Generative AI}{: LangChain, LangGraph, CrewAI, Google Agent Development Kit (ADK), Google Gen AI SDK, retrieval-augmented generation (RAG), EmbeddingGemma embeddings, pgvector, ChromaDB, agentic workflows (planning, memory, tool use), LLM evaluation} \\
     \textbf{Cloud \& Data}{: AWS (S3, EC2, Lambda, SageMaker, Redshift, EMR, ECS/ECR, IAM, CloudWatch), Google Cloud (BigQuery, Managed Service for Apache Spark / Dataproc, Dataflow, Pub/Sub, Cloud Storage, Vertex AI), Apache Spark, Modal} \\
     \textbf{AI Developer Tools}{: Claude Code, GitHub Copilot, Gemini Antigravity, Gemini API} \\
     \textbf{Model Deployment}{: ONNX, NCNN, INT8 quantization, embedded/edge inference} \\
     \textbf{Simulation \& Robotics}{: MuJoCo, Gymnasium, OpenCV, Raspberry Pi} \\
     \textbf{Backend \& Infrastructure}{: Docker, Kubernetes (minikube), FastAPI, Django, Flask, Weights \& Biases, Hydra, Git, Linux, NumPy, pandas}
    }}
 \end{itemize}

%-----------PROJECTS-----------
\section{Projects}
    \resumeSubHeadingListStart

      \resumeProjectHeading
          {\textbf{RoPE-DiT: Diffusion Transformer with 2D RoPE}}{\href{https://github.com/divyanshuklai/Custom2DRoPEDiT}{\uline{github.com/divyanshuklai/Custom2DRoPEDiT}}}
          \resumeItemListStart
            \resumeItem{Developed a 5M-parameter Diffusion Transformer from scratch in PyTorch: patchifier, timestep and class-label embedders, AdaLN-Zero conditioning, and a custom multi-head attention module.}
            \resumeItem{Built a 2D rotary positional embedding that rotates queries and keys by each patch's height and width coordinates, splitting the head dimension into row and column halves so attention sees relative 2D distance.}
            \resumeItem{Trained on CIFAR-10 with a rectified flow objective and sampled with a hand-written Euler solver and classifier-free guidance.}
          \resumeItemListEnd

      \resumeProjectHeading
          {\textbf{Zero-Shot ICE RMA: Quadruped Locomotion RL}}{\href{https://github.com/divyanshuklai/zero-shot-ice-rma}{\uline{github.com/divyanshuklai/zero-shot-ice-rma}}}
          \resumeItemListStart
            \resumeItem{Built a MuJoCo environment for the Unitree Go1 quadruped with 30-dim proprioceptive observations, 12-dim joint-position actions, fractal terrain, and the Rapid Motor Adaptation bioenergetics reward with curriculum-scheduled penalties.}
            \resumeItem{Debugged a PPO baseline that ran 42M steps without learning a gait, tracing the exploding value loss to unnormalized observations and penalty dominance.}
          \resumeItemListEnd

      \resumeProjectHeading
          {\textbf{Fashion Search: Conversational Product Discovery}}{\href{https://github.com/divyanshuklai/fashion_search}{\uline{github.com/divyanshuklai/fashion\_search}}}
          \resumeItemListStart
            \resumeItem{Built a FastAPI chat app answering natural-language product queries across 10 Indian e-commerce stores with Tavily search, Gemini 2.5 Flash extraction, and BeautifulSoup scraping.}
            \resumeItem{Wrote a priority-based image and price extractor that tries JSON-LD product schema, then Open Graph tags, then keyword-matched image tags, as the baseline for an agentic version with planning, memory, and tool use.}
          \resumeItemListEnd

    \resumeSubHeadingListEnd

%-----------ACHIEVEMENTS-----------
\section{Achievements}
  \resumeSubHeadingListStart
    \resumeItem{Won the internal (intra-campus) round of the Smart India Hackathon at Vellore Institute of Technology, 2024.}
  \resumeSubHeadingListEnd

\end{document}
